% Official solution, 2019 DA solutions pp. 8-15; the source's equation
% numbers (2.1)-(2.55) are kept. Photometric phases are written varphi_{1,2}
% and position angles phi_{1,2}, as in the source.

\begin{description}
\item[i)] \hfill
  \begin{description}
  \item[a)] See the following table. \hfill(10 p)

    \begin{center}
    \begin{tabular}{cccccc}
      \hline
      event no. & contacts & components & BJD & $\varphi_1$ & $\varphi_2$ \\
      \hline
      1 & I   & A, B & 2455476.1096 & 0.1226 & 0.9747 \\
        & II  & A, C & 2455476.4245 & 0.4703 & 0.9816 \\
        & III & A, B & 2455477.9677 & 0.1743 & 0.0155 \\
        & IV  & A, B & 2455478.4722 & 0.7313 & 0.0266 \\
      2 & I   & A, B & 2455521.5217 & 0.2643 & 0.9734 \\
      3 & III & A, C & 2455568.9434 & 0.6248 & 0.0163 \\
      4 & I   & A, C & 2455612.4733 & 0.6882 & 0.9736 \\
        & III & A, C & 2455614.3571 & 0.7682 & 0.0150 \\
      5 & III & A, B & 2455659.9241 & 0.0808 & 0.0171 \\
        & IV  & A, C & 2455660.2422 & 0.4320 & 0.0241 \\
      \hline
    \end{tabular}
    \end{center}

  \item[b)] See the following table. \hfill(5 p)

    \begin{center}
    \begin{tabular}{cc}
      \hline
      event no. & closer component \\
      \hline
      1 & A \\
      2 & A \\
      3 & A \\
      4 & A \\
      5 & A \\
      \hline
    \end{tabular}
    \end{center}

    All these events occur close to $\varphi_2 = 0$ (or $\varphi_2 = 1$)
    which means by definition that star A eclipses stars B and C. Therefore,
    star A is closer to the observer.

  \item[c)] In the moments of the 1st and last (4th) contacts:
    \[ R_A + R_{B,C} = d_{\mathrm{A-B,C}},
       \tag*{(2.1)} \]
    \hfill(1 p)\\
    while in the case of the 2nd and 3rd (inner) contacts:
    \[ R_A - R_{B,C} = d_{\mathrm{A-B,C}},
       \tag*{(2.2)} \]
    \hfill(1 p)\\
    where $d_{\mathrm{A-B,C}}$ stands for the sky-projected distance of the
    disks of the occulting star A and occulted component B or C.

    Let $\vec{r}_1$ the radius vector directed from star B toward star C,
    and $\vec{r}_2$ another radius vector directed from the centre of mass
    of stars B and C toward star A. By the use of these two (Jacobian)
    vectors, the position vectors connecting stars B and C with star A, can
    be written as:
    \[ \vec{r}_{\mathrm{BA}} = \vec{r}_A - \vec{r}_B
       = \vec{r}_2 + \frac{m_C}{m_{BC}}\,\vec{r}_1,
       \tag*{(2.3)} \]
    \hfill(1 p)
    \[ \vec{r}_{\mathrm{CA}} = \vec{r}_A - \vec{r}_C
       = \vec{r}_2 - \frac{m_B}{m_{BC}}\,\vec{r}_1.
       \tag*{(2.4)} \]
    \hfill(1 p)

    Using the facts that ($i_{\mathrm{rel}} = 0^\circ$) and, furthermore,
    $i_1 = i_2 = 90^\circ$, it is worthy to introduce the orbital plane as
    reference frame. Let the origin of our frame of reference be the centre
    of mass of the close binary (denoted as $O_2$ in figure above).
    % FIGURE NEEDED: geometry sketch — axes X (toward the observer) and Y in
    % the orbital plane, origin O_2 at the inner-binary centre of mass, stars
    % B and C with Jacobian vector r_1 and angle phi_1, outer centre of mass
    % O_3 with vector r_2 and angle phi_2, vectors r_BA, r_CA to star A;
    % 2019 DA solutions, page 9.
    Furthermore, let axis $X$ directed toward the line of sight. Therefore,
    axis $Y$ is located in the tangential plane of the sky.

    In this frame of reference, the components of vectors $\vec{r}_1$ and
    $\vec{r}_2$ can simply be written as
    \[ \vec{r}_1 = a_1 \left[ \cos(\phi_1); \sin(\phi_1) \right],
       \tag*{(2.5)} \]
    \hfill(1 p)
    \[ \vec{r}_2 = a_2 \left[ \cos(\phi_2); \sin(\phi_2) \right],
       \tag*{(2.6)} \]
    \hfill(1 p)\\
    as $\phi_1$ takes the values of $2k\pi$ in the moments, when star C
    eclipses star B, and similarly, $\phi_2 = 2k\pi$ when star A
    ``eclipses'' the centre of mass of the inner binary.

    One can also notice that there is a simple relation between these
    position angles and the photometric phases defined above as
    \[ \phi_{1,2} = 2\pi\varphi_{1,2}
       \tag*{(2.7)} \]
    \hfill(1 p)

    We are interested in the sky-projected distances of the stellar disks,
    which, in this frame of reference, are equal to the $y$ components of
    the vector equations (2.3 and 2.4). Accordingly,
    \[ d_{\mathrm{A-B}} = \left| a_2 \left[ \sin(2\pi\varphi_2)
       + \frac{m_C}{m_{BC}}\,\frac{a_1}{a_2}\,\sin(2\pi\varphi_1)
       \right] \right|,
       \tag*{(2.8)} \]
    \hfill(1 p)
    \[ d_{\mathrm{A-C}} = \left| a_2 \left[ \sin(2\pi\varphi_2)
       - \frac{m_B}{m_{BC}}\,\frac{a_1}{a_2}\,\sin(2\pi\varphi_1)
       \right] \right|.
       \tag*{(2.9)} \]
    \hfill(1 p)

    \emph{From this point we can use different combinations of eclipsing
    events given in the table. Therefore, here we use one possible order
    only for illustration.}

    We notice that both outer contacts (i.e.\ contacts I and IV) of the
    first event stars A and B are involved. Therefore, we can write the same
    equation for both contacts as follows:
    \[ \frac{R_A + R_B}{a_2} = \left| \sin(2\pi\varphi_2)
       + \frac{m_C}{m_{BC}}\,\frac{a_1}{a_2}\,\sin(2\pi\varphi_1) \right|.
       \tag*{(2.10)} \]
    \hfill(1 p)\\
    Therefore, we have two independent equations and two unknown variables,
    namely
    \[ \frac{R_A + R_B}{a_2}, \quad\text{and}\quad
       \frac{m_C}{m_{BC}}\,\frac{a_1}{a_2} = \frac{a_B}{a_2}.
       \tag*{(2.11)} \]
    \hfill(2 p)

    The computation is as follows: for the 1st contact
    $\sin 2\pi\varphi_2 = -0.158\,296$,
    $\sin 2\pi\varphi_1 = 0.696\,364$, while for the 4th one
    $\sin 2\pi\varphi_2 = 0.166\,356$,
    $\sin 2\pi\varphi_1 = -0.993\,105$. Consequently
    \[ 0.158\,296 - 0.696\,364\,\frac{a_B}{a_2}
       = 0.166\,356 - 0.993\,105\,\frac{a_B}{a_2},
       \tag*{(2.12)} \]
    \hfill(1 p)
    \[ \frac{a_B}{a_2} = 0.027\,161
       \tag*{(2.13)} \]
    \hfill(1 p)\\
    and, therefore,
    \[ \frac{R_A + R_B}{a_2} = 0.1394.
       \tag*{(2.14)} \]
    \hfill(1 p)

    Now, considering the also available 3rd contact moment of the same event
    (in which similarly, stars A and B play the roles), one can get
    \[ \frac{R_A - R_B}{a_2}
       = \left| \sin(2\pi \times 0.0155)
       + 0.027\,161 \times \sin(2\pi \times 0.1743) \right|
       \tag*{(2.15)} \]
    \hfill(1 p)
    \[ \frac{R_A - R_B}{a_2} = 0.1214
       \tag*{(2.16)} \]
    \hfill(1 p)

    Combining these results, we obtain that
    \[ \frac{R_A}{a_2} = \frac{1}{2}\left( \frac{R_A + R_B}{a_2}
       + \frac{R_A - R_B}{a_2} \right) = 0.1304
       \tag*{(2.17)} \]
    \hfill(2 p)
    \[ \frac{R_B}{a_2} = \frac{1}{2}\left( \frac{R_A + R_B}{a_2}
       - \frac{R_A - R_B}{a_2} \right) = 0.0090
       \tag*{(2.18)} \]
    \hfill(2 p)

    We repeat a very similar calculation for the 3rd contacts of events 3
    and 4, in which cases star B is substituted with star C. For the 3rd
    contact of the third event $\sin 2\pi\varphi_2 = 0.102\,237$,
    $\sin 2\pi\varphi_1 = -0.706\,218$, while the 3rd contact of the fourth
    event $\sin 2\pi\varphi_2 = 0.094\,108$,
    $\sin 2\pi\varphi_1 = -0.993\,469$. Therefore,
    \[ 0.102\,237 + 0.706\,218\,\frac{a_C}{a_2}
       = 0.094\,108 + 0.993\,469\,\frac{a_C}{a_2}
       \tag*{(2.19)} \]
    \hfill(1 p)
    \[ \frac{a_C}{a_2} = 0.028\,298
       \tag*{(2.20)} \]
    \hfill(1 p)\\
    and, accordingly,
    \[ \frac{R_A - R_C}{a_2} = 0.1222
       \tag*{(2.21)} \]
    \hfill(1 p)

    Taking into account that $R_A/a_2$ is already known from the previous
    stage, the dimensionless relative radius of star C can be obtained
    simply as
    \[ \frac{R_C}{a_2} = \frac{R_A}{a_2} - \frac{R_A - R_C}{a_2}
       = 0.1304 - 0.1222 = 0.0082
       \tag*{(2.22)} \]
    \hfill(2 p)

    The other possibility is, however, that e.g.\ from the 1st contact of
    event 4 we calculate the sum $(R_A + R_C)/a_2$, as
    \[ \frac{R_A + R_C}{a_2}
       = \left| \sin(2\pi \times 0.9736)
       - 0.028\,298 \times \sin(2\pi \times 0.6882) \right|
       \tag*{(2.23)} \]
    \[ \frac{R_A + R_C}{a_2} = 0.1389,
       \tag*{(2.24)} \]
    and then we obtain that
    \[ \frac{R_A}{a_2} = \frac{1}{2}\left( \frac{R_A + R_C}{a_2}
       + \frac{R_A - R_C}{a_2} \right) = 0.1306,
       \tag*{(2.25)} \]
    \[ \frac{R_C}{a_2} = \frac{1}{2}\left( \frac{R_A + R_C}{a_2}
       - \frac{R_A - R_C}{a_2} \right) = 0.0084.
       \tag*{(2.26)} \]

    \begin{quote}
    \emph{Several additional, equivalently acceptable scenarios.}

    We can write three equations for $(R_A + R_B)/a_2$:
    \begin{align*}
      \frac{R_A + R_B}{a_2}
        &= 0.158\,296 - 0.696\,364\,\frac{a_B}{a_2}
        && \text{(1st event Ist contact)} \\
      \frac{R_A + R_B}{a_2}
        &= 0.166\,356 - 0.993\,105\,\frac{a_B}{a_2}
        && \text{(1IV)} \\
      \frac{R_A + R_B}{a_2}
        &= 0.166\,433 - 0.995\,966\,\frac{a_B}{a_2}
        && \text{(2I)}
    \end{align*}
    Combining them one can get:
    \begin{center}
    \begin{tabular}{ccc}
      \hline
      contacts & $a_B/a_2$ & $(R_A + R_B)/a_2$ \\
      \hline
      1I--1IV & 0.027\,161 & 0.1394 \\
      1I--2I  & 0.027\,159 & 0.1394 \\
      1IV--2I & 0.026\,988 & 0.1396 \\
      \hline
    \end{tabular}
    \end{center}
    Similarly, for $(R_A - R_B)/a_2$:
    \begin{align*}
      \frac{R_A - R_B}{a_2}
        &= 0.097\,235 + 0.889\,001\,\frac{a_B}{a_2}
        && \text{(1III)} \\
      \frac{R_A - R_B}{a_2}
        &= 0.107\,236 + 0.486\,152\,\frac{a_B}{a_2}
        && \text{(5III)}
    \end{align*}
    Combining them one can get:
    \begin{center}
    \begin{tabular}{ccc}
      \hline
      contacts & $a_B/a_2$ & $(R_A - R_B)/a_2$ \\
      \hline
      1III--5III & 0.024\,826 & 0.1193 \\
      \hline
    \end{tabular}
    \end{center}
    Now for $(R_A + R_C)/a_2$:
    \begin{align*}
      \frac{R_A + R_C}{a_2}
        &= 0.165\,116 - 0.925\,553\,\frac{a_C}{a_2}
        && \text{(4I)} \\
      \frac{R_A + R_C}{a_2}
        &= 0.150\,847 - 0.414\,376\,\frac{a_C}{a_2}
        && \text{(5IV)}
    \end{align*}
    The only combination gives:
    \begin{center}
    \begin{tabular}{ccc}
      \hline
      contacts & $a_C/a_2$ & $(R_A + R_C)/a_2$ \\
      \hline
      4I--5IV & 0.027\,914 & 0.1393 \\
      \hline
    \end{tabular}
    \end{center}
    Finally, for $(R_A - R_C)/a_2$:
    \begin{align*}
      \frac{R_A - R_C}{a_2}
        &= 0.115\,353 + 0.185\,529\,\frac{a_C}{a_2}
        && \text{(1II)} \\
      \frac{R_A - R_C}{a_2}
        &= 0.102\,237 + 0.706\,218\,\frac{a_C}{a_2}
        && \text{(3III)} \\
      \frac{R_A - R_C}{a_2}
        &= 0.094\,108 + 0.993\,469\,\frac{a_C}{a_2}
        && \text{(4III)}
    \end{align*}
    The combinations are:
    \begin{center}
    \begin{tabular}{ccc}
      \hline
      contacts & $a_C/a_2$ & $(R_A - R_C)/a_2$ \\
      \hline
      1II--3III  & 0.025\,190 & 0.1200 \\
      1II--4III  & 0.026\,295 & 0.1202 \\
      3III--4III & 0.028\,299 & 0.1222 \\
      \hline
    \end{tabular}
    \end{center}
    \emph{End additional, equivalently acceptable scenarios.}
    \end{quote}

    Now, we are in the position to calculate the inner mass ratio $q_1$, as
    follows:
    \[ q_1 = \frac{a_B/a_2}{a_C/a_2}
       = \frac{0.027\,161}{0.028\,298} = 0.9598
       \tag*{(2.27)} \]
    \hfill(2 p)

    Furthermore, we can also get the ratio of the semi-major axes as
    \[ \frac{a_1}{a_2} = \frac{a_B}{a_2} + \frac{a_C}{a_2}
       = 0.027\,161 + 0.028\,298 = 0.055\,459
       \tag*{(2.28)} \]
    \hfill(2 p)

    The requested results with the range of the full marks (summary):
    \begin{align*}
      0.1150 &\le \frac{R_A}{a_2} \le 0.1450 \\
      0.0075 &\le \frac{R_B}{a_2} \le 0.0105 \\
      0.0070 &\le \frac{R_C}{a_2} \le 0.0100 \\
      0.045 &\le \frac{a_1}{a_2} \le 0.065 \\
      0.80 &\le q_1 \le 1.25
    \end{align*}

  \item[d)] The quickest method to obtain $q_2$ comes with the double use of
    Kepler's third law. Writing this law for both the inner and outer
    orbits, and dividing them one can get, that
    \[ \left( \frac{a_1}{a_2} \right)^3 \left( \frac{P_2}{P_1} \right)^2
       = \frac{m_{BC}}{m_{ABC}} = \frac{q_2}{1 + q_2}.
       \tag*{(2.29)} \]
    \hfill(4 p)\\
    Therefore
    \[ q_2 = \frac{\left( \frac{a_1}{a_2} \right)^3
       \left( \frac{P_2}{P_1} \right)^2}
       {1 - \left( \frac{a_1}{a_2} \right)^3
       \left( \frac{P_2}{P_1} \right)^2}
       \tag*{(2.30)} \]
    \hfill(2 p)
    \[ q_2 \approx \frac{0.055\,459^3 \times 50.206\,763^2}
       {1 - 0.055\,459^3 \times 50.206\,763^2} \approx 0.7543
       \tag*{(2.31)} \]
    \hfill(2 p)
  \end{description}

\item[ii)] The circular velocity of star A is
  \[ v_A = \frac{2\pi}{P_2}\,r_A.
     \tag*{(2.32)} \]
  \hfill(1 p)\\
  The amplitude of the RV curve is equal to maximum value of the
  line-of-sight component of that velocity, i.e.
  \[ K_A = v_A \sin i_2.
     \tag*{(2.33)} \]
  \hfill(1 p)

  The sinusoidal component in the occurrence of the eclipsing minima times
  comes from the light-travel time effect caused by the revolution of stars
  B and C around the centre of mass of the whole triple system. In this
  regard the motion of components B and C can simply be substituted by the
  movement of their centre of mass along the outer orbit. Therefore, the
  total variation of the inner binary's distance to the Earth is the
  line-of-sight component of the diameter of the orbit of the centre of mass
  of the inner binary around the centre of mass of the complete triple
  system, i.e.
  \[ \Delta z_{BC} = 2 r_{BC} \sin i_2.
     \tag*{(2.34)} \]
  \hfill(2 p)\\
  Therefore, the amplitude of the light-travel time sine wave is
  \[ A_{\mathrm{ETV}} = \frac{\Delta z}{2c}
     = \frac{r_{BC} \sin i_2}{c}
     \;\rightarrow\;
     r_{BC} = \frac{c A_{\mathrm{ETV}}}{\sin i_2}.
     \tag*{(2.35)} \]
  \hfill(2 p)

  Such a way, using that
  \[ q_2 = \frac{m_{BC}}{m_A} = \frac{r_A}{r_{BC}},
     \tag*{(2.36)} \]
  \hfill(1 p)\\
  one can obtain, that
  \[ q_2 = \frac{P_2 K_A}{2\pi c A_{\mathrm{ETV}}}.
     \tag*{(2.37)} \]
  \hfill(1 p)\\
  which gives a second, independent determination of the outer mass ratio
  $q_2$.

  The uncertainty can be estimated either as
  \[ \Delta q_2 = q_2 \sqrt{
       \left( \frac{\Delta P_2}{P_2} \right)^2
       + \left( \frac{\Delta K_A}{K_A} \right)^2
       + \left( \frac{\Delta A_{\mathrm{ETV}}}{A_{\mathrm{ETV}}} \right)^2 }
     \tag*{(2.38)} \]
  or
  \[ \Delta q_2 = q_2 \left(
       \left| \frac{\Delta P_2}{P_2} \right|
       + \left| \frac{\Delta K_A}{K_A} \right|
       + \left| \frac{\Delta A_{\mathrm{ETV}}}{A_{\mathrm{ETV}}} \right|
       \right)
     \tag*{(2.39)} \]
  \hfill(2 p)

  With numerical values:
  \[ q_2 = 0.621 \pm 0.047 \quad\text{or}\quad q_2 = 0.621 \pm 0.048 \]
  \hfill(1 p)

  Then, the semi-major axis of the outer orbit can be calculated in the
  following alternative ways:
  \[ a_2 = r_A\,\frac{1 + q_2}{q_2}
     = \frac{P_2}{2\pi}\,\frac{K_A}{\sin i_2}\,\frac{1 + q_2}{q_2}
     \tag*{(2.40)} \]
  \[ a_2 = r_{BC}\,(1 + q_2)
     = \frac{c A_{\mathrm{ETV}}}{\sin i_2}\,(1 + q_2)
     \tag*{(2.41)} \]
  \[ a_2 = r_A + r_{BC}
     = \frac{P_2}{2\pi}\,\frac{K_A}{\sin i_2}
     + \frac{c A_{\mathrm{ETV}}}{\sin i_2},
     \tag*{(2.42)} \]
  One of the three equations above: \hfill(2 p)

  We assumed that $i_2 = 90^\circ$, therefore $\sin i_2 = 1$. Then
  \[ a_2 = (60.712 \pm 2.851) \times 10^6\ \mathrm{km}
     = (87.293 \pm 4.099)\ R_\odot
     = (0.406 \pm 0.019)\ \mathrm{AU} \]
  \hfill(2 p)

  The uncertainties from the three different formulae above:
  \[ \Delta a_2 = a_2 \sqrt{
       \left( \frac{\Delta P_2}{P_2} \right)^2
       + \left( \frac{\Delta K_A}{K_A} \right)^2
       + \left( \frac{\Delta q_2}{q_2 (1 + q_2)} \right)^2 }
     \tag*{(2.43)} \]
  \[ \Delta a_2 \approx 2.851 \times 10^6\ \mathrm{km}
     \approx 4.10\ R_\odot \approx 0.019\ \mathrm{AU} \]
  \[ \Delta a_2 = a_2 \sqrt{
       \left( \frac{\Delta A_{\mathrm{ETV}}}{A_{\mathrm{ETV}}} \right)^2
       + \left( \frac{\Delta q_2}{1 + q_2} \right)^2 }
     \tag*{(2.44)} \]
  \[ \Delta a_2 \approx 4.946 \times 10^6\ \mathrm{km}
     \approx 7.11\ R_\odot \approx 0.033\ \mathrm{AU} \]
  \[ \Delta a_2 = \sqrt{
       \left( \frac{\Delta P_2}{2\pi}\,K_A \right)^2
       + \left( \frac{P_2}{2\pi}\,\Delta K_A \right)^2
       + \left( c\,\Delta A_{\mathrm{ETV}} \right)^2 }
     \tag*{(2.45)} \]
  \[ \Delta a_2 \approx 2.849 \times 10^6\ \mathrm{km}
     \approx 4.10\ R_\odot \approx 0.019\ \mathrm{AU} \]
  One of the three formulae above with its uncertainty: \hfill(3 p)

  In what follows, we use the smallest uncertainties (first row above) but
  the others, and also the nonquadratic ones are acceptable, too.

  The total mass of the triple now can be calculated directly from Kepler's
  third law or, as an alternative way, one can obtain it from the equations
  on the centripetal accelerations.
  \[ r_A\,\frac{4\pi^2}{P_2^2} = \frac{G m_{BC}}{a_2^2}
     \tag*{(2.46)} \]
  \[ r_{BC}\,\frac{4\pi^2}{P_2^2} = \frac{G m_A}{a_2^2},
     \tag*{(2.47)} \]
  where the only unknown quantities are the masses, so they can be
  calculated separately from the two equations or, summing the masses we can
  get back (formally) Kepler's third law:
  \[ m_{ABC} = \frac{4\pi^2}{G}\,\frac{a_2^3}{P_2^2}
     = (4.312 \pm 0.607)\ M_\odot,
     \tag*{(2.48)} \]
  \hfill(2 p)\\
  from which
  \[ m_A = \frac{m_{ABC}}{1 + q_2} = (2.660 \pm 0.383)\ M_\odot
     \tag*{(2.49)} \]
  \hfill(1 p)
  \[ m_{BC} = m_{ABC}\,\frac{q_2}{1 + q_2} = (1.652 \pm 0.245)\ M_\odot
     \tag*{(2.50)} \]
  \hfill(1 p)

\item[iii)] We obtained the total mass ($m_{BC}$) of the inner binary in the
  second part of the problem. Combining this with the inner mass ratio
  ($q_1$) obtained in the first part, one can get that
  \[ m_B = \frac{m_{BC}}{1 + q_1}
     = \frac{1.652\ M_\odot}{1 + 0.960} = 0.843\ M_\odot
     \tag*{(2.51)} \]
  \hfill(3 p)\\
  and
  \[ m_C = q_1 \times m_B = 0.960 \times 0.843\ M_\odot = 0.809\ M_\odot
     \tag*{(2.52)} \]
  \hfill(3 p)

  The dimensionless radius of each star relative to the semi-major axis was
  calculated in the first part, while the semi-major axis of the outer orbit
  ($a_2$) was obtained in the second part. Their multiplication gives the
  physical radii of the stars:
  \[ R_A = 0.1304 \times 87.29\ R_\odot = 11.381\ R_\odot
     \tag*{(2.53)} \]
  \hfill(3 p)
  \[ R_B = 0.0090 \times 87.29\ R_\odot = 0.786\ R_\odot
     \tag*{(2.54)} \]
  \hfill(3 p)
  \[ R_C = 0.0082 \times 87.29\ R_\odot = 0.712\ R_\odot
     \tag*{(2.55)} \]
  \hfill(3 p)
\end{description}
